% !TEX root = ../thesis.tex
% ==============================================================================
%  APPENDIX: SENTENCE-LEVEL ADAPTATION OF LIME (appendices/appendix_lime_adaptation.tex)
% ==============================================================================
%  This appendix details the mathematical formulation and implementation
%  for adapting LIME from token perturbations to sentence-level attribution
%  under black-box generative constraints.
% ==============================================================================

\section{تطبیق \lr{LIME} به سطح جمله و دسترسی مسدود}\label{app:appendix_lime_adaptation}

همان‌طور که در فصل‌های~\ref{chap:chapter3_literature} و~\ref{chap:chapter4_methodology} استدلال گردید، به دلیل نارسایی‌های معناییِ رویکردهای توکن‌محور، ما فضای مسئله و واحد تحلیل را در این پژوهش به سطح جمله (مفاهیم گزاره‌ای) ارتقا دادیم. با توجه به این تغییر ساختاری، مقایسه‌ی مستقیم مفسر جمله‌محور ما با نسخه‌ی خام و توکن‌محور الگوریتم \lr{LIME}~\cite{ribeiro2016lime} از اساس نامعتبر و غیرعادلانه خواهد بود؛ چرا که \lr{LIME} کلاسیک صرفاً تأثیر رخداد توکن‌های منفرد و نوسانات آماری آن‌ها را می‌سنجد و از ارزیابی تأثیر حضور یا غیاب یک مفهوم یکپارچه‌ی معنایی عاجز است.

لذا برای فراهم آوردن یک خط‌پایه‌ی مقایسه‌ای معنادار، ما وظیفه‌ی این الگوریتم را با صورت‌مسئله‌ی خویش تطبیق داده و واحد استخراج آن را از توکن به جمله ارتقا دادیم. این تطبیق باعث می‌شود که انتخاب‌های \lr{LIME} از سطح نوسانات پرنویز نحوی فراتر رفته و به مفاهیم منطقی قابل‌درک برای انسان تبدیل شوند. این هم‌سان‌سازی در واحد تحلیل، امکان یک سنجش علمی و هم‌تراز را میان استراتژی تقریب محلی (\lr{LIME}) و تقریب سراسری (\lr{ESCI}) فراهم می‌آورد. در ادامه، فرمول‌بندی ریاضی این ارتقا و نحوه‌ی پیاده‌سازی آن در دو سناریوی جعبه‌خاکستری و جعبه‌سیاهِ مسدود تشریح می‌شود.


\subsection{مبانی ریاضی \lr{LIME} و ارتقا به سطح جمله}\label{app:lime_formulation}

الگوریتم \lr{LIME} بر پایه‌ی تقریب رفتار یک مدل پیچیده ($f_\theta$) در همسایگی محلیِ یک ورودی مشخص ($\mathbf{x}$) به ازای تولید یک گزاره‌ی هدف معین ($y_t$) از طریق یک مدل جانشین ذاتاً تفسیرپذیر ($g \in G$، غالباً رگرسیون خطی خلوت) استوار است.

مطابق با تعاریف ارائه‌شده در بخش~\ref{sec:1-1-1} (تعریف~\ref{def:chap1_proposition_space} و معادله‌ی~\ref{eq:chap1_prompt_decomp}) و خط‌لوله‌ی تقطیع در بخش~\ref{sec:4-2-1}، ابتدا پرامپت ورودی کاربر به جای توالی خرد کلمات، متشکل از دنباله‌ی توکن‌های $\mathbf{x} = (u_1, u_2, \dots, u_{T_{\text{\lr{in}}}}) \in \mathcal{V}^{T_{\text{\lr{in}}}}$، پس از اعمال عملگر تقطیع نحوی $\Pi$، به صورت دنباله‌ای از $N$ جمله تعریف می‌شود:
\begin{equation}\label{eq:app_prompt_decomp}
    \mathbf{x} = (x_1, x_2, \dots, x_N) \in \mathcal{S}^N
\end{equation}
که در آن هر جمله‌ی ورودی $x_i \in \mathcal{S}$ خود زیردنباله‌ای پیوسته از توکن‌های ورودی به فرمی به صورت $x_i = (u_{\tau_{i-1}+1}, \dots, u_{\tau_i})$ با شاخص‌های $0 = \tau_0 < \tau_1 < \cdots < \tau_N = T_{\text{\lr{in}}}$ و طول $|x_i| = \tau_i - \tau_{i-1}$ است. همچنین بر اساس تعریف~\ref{def:chap1_proposition_space} (معادله‌ی~\ref{eq:chap1_response_decomp}) و تعریف~\ref{def:chap1_target_subset}، پاسخ تولیدشده توسط مدل به صورت دنباله‌ای از $M$ جمله $\mathbf{y} = (y_1, \dots, y_M) \in \mathcal{S}^M$ افراز شده که در آن هر جمله به فرم مذکور $y_j = (v_{\sigma_{j-1}+1}, \dots, v_{\sigma_j})$ با شاخص‌های مرزی $0 = \sigma_0 < \cdots < \sigma_M = T_{\text{\lr{out}}}$ تعریف می‌شود، و گزاره‌ی هدف خروجی که تحلیل‌گر قصد تبیین منشأ آن را دارد، با $y_t \in \mathcal{S}$ ($t \in T \subseteq \{1, \dots, M\}$) و توکن‌های $(v_{\sigma_{t-1}+1}, \dots, v_{\sigma_t})$ مشخص می‌گردد.

در این تعریف، هر ویژگی در رگرسیون موضعی ما، معادل با یک جمله‌ی کامل $x_i$ است. برای سنجش اهمیت این ویژگی‌ها، الگوریتم \lr{LIME} نیازمند نمونه‌برداری و ایجاد اختلال در ورودی است تا تأثیر آن در خروجی مدل را مشاهده کند. بدین منظور، مطابق با فضای فرضیات انتساب در تعریف~\ref{def:chap1_mask_space} (یعنی $\mathcal{M} = {\{0, 1\}}^N$)، مجموعه‌ای از $K$ بردار ماسک دودویی $\mathbf{m}^{(k)} = {(m_1^{(k)}, m_2^{(k)}, \dots, m_N^{(k)})}^\top \in {\{0, 1\}}^N$ تولید می‌گردد که در آن‌ها، درایه‌ی $m_i^{(k)} = 1$ به معنای حفظ جمله‌ی $x_i$ و درایه‌ی $m_i^{(k)} = 0$ به معنای حذف فیزیکی آن از پرامپت است.

با توجه به ارتقای واحد تحلیل به سطح جمله، طول بردارهای ویژگی (میانگین $N \approx 15$) بسیار کوتاه‌تر از فضای توکن‌هاست. با این حال، فضای ترکیبیاتی $\mathcal{O}(2^N)$ همچنان می‌تواند بزرگ باشد. به منظور مدیریت هزینه‌های محاسباتیِ استنتاج از مدل‌های زبانی بزرگ و در عین حال تضمین ثبات آماریِ مدل رگرسیون (تأمین درجه‌ی آزادی کافی)، تعداد نمونه‌ها به صورت $K = \min(2^N, 1000)$ تنظیم گردید. با اعمال هر ماسک، یک پرامپت مخدوش‌شده به فرم متنی زیر ساخته می‌شود که منطبق بر زیرمجموعه‌ی پیش‌آیند ورودی $\mathbf{x}_S^{(k)} = \mathbf{x} \odot \mathbf{m}^{(k)}$ در معادله‌ی~\ref{eq:chap1_subset_mask} است:
\begin{equation}\label{eq:app_lime_perturbed_text}
    \tilde{\mathbf{x}}^{(k)} = \text{\lr{Concat}}\left( \{x_i \in \mathbf{x} \mid m_i^{(k)} = 1, \, \forall i \in \{1, \dots, N\} \} \right)
\end{equation}

بر اساس تئوری اصلی \lr{LIME}، مدل خطی باید بر روی نمونه‌هایی که به داده‌ی اصلی شبیه‌ترند، تمرکز بیشتری داشته باشد. این وزن‌دهی محلی ($\pi_{\mathbf{x}}$) نه بر روی متن، بلکه مستقیماً در فضای گسسته‌ی ماسک‌های دودویی انجام می‌پذیرد. اگر $\mathbf{1} = {(1, 1, \dots, 1)}^\top \in {\{1\}}^N$ نشان‌دهنده‌ی بردار ماسک اصلی (بدون هیچ حذفیات، معادل پرامپت کامل) باشد، فاصله‌ی نمونه‌ی $k$-ام با استفاده از معیار فاصله‌ی کسینوسی محاسبه می‌گردد:
\begin{equation}\label{eq:app_lime_distance}
    d_k = 1 - \frac{{\mathbf{m}^{(k)}}^\top \mathbf{1}}{\|\mathbf{m}^{(k)}\|_2 \|\mathbf{1}\|_2}
\end{equation}
از آنجایی که این فاصله ذاتاً در بازه‌ی $[0, 1]$ نرمال‌سازی شده است، وزن هر نمونه ($w_k$) با تبعیت از تنظیمات استانداردِ \lr{LIME} برای فواصل نرمال، با استفاده از یک هسته‌ی نمایی\enfootnote{exponential kernel} با پهنای ثابت به دست می‌آید:
\begin{equation}\label{eq:app_lime_kernel}
    w_k = \exp\left(-\frac{d_k^2}{\sigma_{\text{\lr{kernel}}}^2}\right), \quad \text{حيث} \quad \sigma_{\text{\lr{kernel}}} = 0.25
\end{equation}

در نهایت، با استخراج یک امتیاز هدف عددی ($r_k \in \mathbb{R}$) که نشان‌دهنده‌ی واکنش مدل به پرامپت مخدوش‌شده‌ی $\tilde{\mathbf{x}}^{(k)}$ به ازای گزاره‌ی هدف $y_t$ است، یک مدل رگرسیون لاسو آموزش می‌بیند. این مدل با کمینه‌سازی تابع زیان زیر، بردار ضرایب $\boldsymbol{\beta} = {(\beta_1, \beta_2, \dots, \beta_N)}^\top \in \mathbb{R}^N$ (که همان نمرات تخصیص و اهمیت جملات، $s_i = \beta_i$، هستند) را تخمین می‌زند:
\begin{equation}\label{eq:app_lime_lasso}
    \min_{\boldsymbol{\beta} \in \mathbb{R}^N} \sum_{k=1}^K w_k {\left( r_k - {\mathbf{m}^{(k)}}^\top \boldsymbol{\beta} \right)}^2 + \lambda \|\boldsymbol{\beta}\|_1
\end{equation}
در این معادله، $r_k$ متغیر وابسته (امتیاز واکنش مدل به پرامپت مخدوش به ازای خروجی هدف)، $\mathbf{m}^{(k)}$ بردار ماسک دودویی (به عنوان متغیرهای مستقل ورودی)، و $\boldsymbol{\beta}$ بردار ضرایب رگرسیون است که مقادیر نهایی اهمیت (نمرات تخصیص) برای هر جمله‌ی پرامپت را نشان می‌دهد. همچنین، $\lambda$ ابرپارامتر منظم‌ساز است که میزان خلوت‌سازی (صفر شدن ضرایب برای جملات بی‌اهمیت) را کنترل می‌کند.

برای یافتن بهینه‌ترین مقدار $\lambda$ که از بیش‌برازش جلوگیری کند، از الگوریتم \lr{LassoCV} (رگرسیون لاسو مجهز به اعتبارسنجی متقاطع) استفاده شده است. در این روش، داده‌های نمونه‌برداری‌شده به ۵ بخش مساوی (فولد) تقسیم می‌شوند؛ در هر مرحله، مدل روی ۴ بخش آموزش دیده و روی بخش پنجم ارزیابی می‌گردد و این چرخه برای تمام بخش‌ها تکرار می‌شود تا در نهایت $\lambda$ای با کمترین خطای تعمیم‌پذیری انتخاب شود.

مهم‌ترین مسئله، نحوه‌ی استخراج متغیر وابسته ($r_k$) است که مرز تمایز میان سناریوی جعبه‌خاکستری (مبتنی بر لوجیت) و جعبه‌سیاه (مبتنی بر خروجی متنی) را تعیین می‌کند. در ادامه، دو رویکرد تعریف مسئله و راه‌حل هر یک را در پیاده‌سازی خود شرح می‌دهیم.


\subsection{\lr{LIME} جعبه‌خاکستری}\label{app:lime_graybox}

در ارزیابی مدل‌های متن‌باز (مانند \lr{Llama-3})، ما به توزیع احتمالات درونی مدل (لوجیت‌های لایه‌ی آخر) دسترسی داریم. این سناریو به عنوان حالت جعبه‌خاکستری تعریف می‌شود. در اینجا، ما می‌خواهیم تأثیر حذف جملات پرامپت را بر روی احتمال تولید دقیق یک جمله‌ی هدف خاص ($y_t$) بسنجیم.

مطابق با تعاریف~\ref{def:chap1_proposition_space} و~\ref{def:chap1_target_subset} در بخش~\ref{sec:1-1-1}، فرض کنید گزاره‌ی هدف خروجی تک‌جمله‌ای $y_t \in \mathcal{S}$ با زیردنباله‌ای از $L_t = \sigma_t - \sigma_{t-1}$ توکن به صورت $y_t = (v_{\sigma_{t-1}+1}, \dots, v_{\sigma_t})$ مشخص شده باشد. متغیر هدف ($r_k$) در رگرسیون معادله‌ی~\ref{eq:app_lime_lasso}، لگاریتم احتمال تولید خودهمبسته‌ی این دنباله، مشروط بر دریافت پرامپت مخدوش‌شده‌ی $\tilde{\mathbf{x}}^{(k)}$ است. برای محاسبه‌ی کارآمد این متغیر و پرهیز از هزینه‌ی تولید توکن‌به‌توکن، ما از تکنیک هدایت معلم‌وار یکپارچه و عبور پیش‌رو استفاده کردیم. در این روش، توکن‌های پرامپت مخدوش و توکن‌های جمله‌ی هدف به یکدیگر متصل شده و به صورت یکپارچه ($\mathbf{I}^{(k)} = [\tilde{\mathbf{x}}^{(k)} \oplus y_t]$) به مدل زبانی تغذیه می‌گردند. در شبکه‌های خودهمبسته، احتمال تولید توکن هدف $v_j$ توسط حالت پنهانِ توکن ماقبل آن پیش‌بینی می‌شود. سپس، لوجیت‌های متناظر با جایگاه توکن‌های هدف استخراج می‌شوند. در نهایت، متغیر $r_k$ مستقیماً با محاسبه‌ی منفی خطای آنتروپی متقاطع (\lr{CE}) بر روی این لوجیت‌ها و جمع زدن لگاریتم احتمالات شرطی تک‌تک توکن‌ها به دست می‌آید:
\begin{equation}\label{eq:app_lime_graybox_target}
\begin{split}
    r_k &= \log P_\theta\left(y_t \mid \tilde{\mathbf{x}}^{(k)}\right) \\
    &= \sum_{j=1}^{L_t} \log P_\theta\left(v_{\sigma_{t-1}+j} \mid \tilde{\mathbf{x}}^{(k)}, v_{\sigma_{t-1}+<j}\right) \\
    &= -\sum_{j=1}^{L_t} \text{\lr{CE}}\left(\mathbf{z}_{T_{\text{\lr{in}}}^{(k)}+j-1}, v_{\sigma_{t-1}+j}\right)
\end{split}
\end{equation}
که در آن $T_{\text{\lr{in}}}^{(k)}$ طول توکن‌های پرامپت مخدوش $\tilde{\mathbf{x}}^{(k)}$، و $\mathbf{z}$ بردار لوجیت‌های خروجی است.

همان‌طور که از این معادله پیداست، اگرچه ما در این تطبیق، ورودی را در سطح جمله (به عنوان یک واحد کلان) مخدوش می‌کنیم، اما خروجی و متغیر هدف همچنان از طریق جمع زدنِ نوسانات احتمالات خرد و توکن‌به‌توکن ارزیابی می‌شود. این دوگانگی ساختاری نشان می‌دهد که نسخه‌ی جعبه‌خاکستری \lr{LIME}، علی‌رغم ارتقای ورودی، در حقیقت یک ارزیاب کاملاً جمله‌محور و مفهومی نیست.

باید توجه داشت که اگرچه محاسبه‌ی احتمالات توکن‌به‌توکن از نظر محاسباتی دقیق‌ترین رویکرد برای سنجش نوسانات یک مدل خودهمبسته است، اما این رویکرد به هیچ وجه استاندارد طلایی بلامنازع در وظیفه‌ی تفسیر مفهومی به شمار نمی‌رود. متدولوژی‌هایی که به رخداد دقیق توکن‌ها اتکا می‌کنند، در برابر پدیده‌ی پارافریز آسیب‌پذیرند؛ به این معنا که مدل زبانی ممکن است در اثر تغییر پرامپت، همان مفهوم منطقی را با دایره‌ی واژگان متفاوتی بازتولید کند، اما این روش با مشاهده‌ی کاهش احتمالِ توکن‌های اولیه (به دلیل شیفت توزیع ناشی از کلمات جدید)، آن را به اشتباه به عنوان از بین رفتن مفهوم هدف تلقی نماید. با این وجود، این روش به عنوان کران بالای محاسباتی در رویکردهای اختلال محلی در نظر گرفته می‌شود و انتظار می‌رود در معیارهای ارزیابی مبتنی بر تغییر رخداد توکن‌ها، عملکرد بهینه‌ای داشته باشد.


\subsection{\lr{LIME} جعبه‌سیاه}\label{app:lime_blackbox}

در مدل‌های تجاری و کاملاً مسدود (نظیر \lr{GPT-4o-Mini})، استخراج لوجیت‌ها ناممکن است. در این سناریو، ما تنها می‌توانیم پرامپت مخدوش ($\tilde{\mathbf{x}}^{(k)}$) را از طریق رابط برنامه‌نویسی به مدل ارسال کرده، یک متن خروجی کاملاً جدید ($\tilde{\mathbf{y}}^{(k)} = f_\theta(\tilde{\mathbf{x}}^{(k)})$) دریافت نماییم و سپس بررسی کنیم که این متن جدید تا چه حد بازتولیدکننده‌ی گزاره‌ی هدف اولیه‌ی ما ($y_t$) است.

در این مرحله، انتخاب تابع شباهت برای محاسبه‌ی متغیر $r_k$ چالش تئوریک اصلی است. در ابتدا، استفاده از شباهت کسینوسی میان بردار تعبیه‌ی متراکم گزاره‌ی هدف اصلی ($y_t$) و بردار تعبیه‌ی کل متن جدیدِ تولیدشده ($\tilde{\mathbf{y}}^{(k)}$) یا تک‌تک جملاتش (نظیر استفاده از خروجی‌های مدل‌های \lr{SentenceTransformers}~\cite{reimers2019sentencebert}) مورد آزمون قرار گرفت، اما این رویکرد به دو دلیل ماهوی از خط‌لوله‌ی \lr{LIME} کنار گذاشته شد:
نخست، به دلیل تضاد با ذات الگوریتم؛ چرا که الگوریتم \lr{LIME} از پایه برای یافتن هم‌پوشانی‌های موضعی و واژگانی توسعه یافته است و استفاده از بردارهای متراکم برای کل متن که مفاهیم را به‌شدت انتزاع می‌کنند، ذات و رفتار رگرسیون محلی این الگوریتم را از ریشه‌ی خود منحرف می‌سازد.
دلیل دوم، مسئله‌ی نشت داده و برتری دُوری است؛ از آنجا که در فصل آزمایش‌ها برای سنجش نهایی وفاداری (از طریق محاسبه‌ی جامعیت و کفایت) از معیار شباهت کسینوسی استفاده نموده‌ایم، به‌کارگیری همین تابع به عنوان هدف بهینه‌سازیِ \lr{LIME} در فاز استخراج نمرات، یک برتری ناعادلانه و چرخه‌ای برای این خط پایه ایجاد می‌کرد که اعتبار علمی مقایسه‌ها را مخدوش می‌ساخت.

برای پرهیز از این سوگیری‌ها و در عین حال حفظ دقت الگوریتم در سطح واژگان، ما از معیار \lr{BERTScore}~\cite{zhang2020bertscore} استفاده کردیم. برخلاف معیارهای سنتی که صرفاً بر تطابق دقیق و دودویی کلمات استوارند، \lr{BERTScore} از تعبیه‌های بافتاری\enfootnote{contextual embeddings} استخراج‌شده از مدل‌های زبانی از‌پیش‌آموزش‌دیده (مانند \lr{RoBERTa}~\cite{liu2021roberta}) بهره می‌برد. در این راستا، رویکردِ بُعد پردازش توکن خود را به گونه‌ای حفظ می‌کند. مکانیسم عمل این معیار بدین صورت است که ابتدا برای تک‌تک توکن‌های موجود در متن تولیدشده ($\tilde{\mathbf{y}}^{(k)}$) و تک‌تک توکن‌های گزاره‌ی هدف ($y_t$)، بردارهای متراکم بافتاری تولید می‌شود. سپس، یک ماتریس شباهت کسینوسی دوبه‌دو میان تمام توکن‌های متن تولیدشده و توکن‌های گزاره‌ی هدف محاسبه می‌گردد. در گام بعد، از طریق یک الگوریتم تطبیق حریصانه\enfootnote{greedy matching}، برای هر توکن در یک متن، شبیه‌ترین توکن در متن مقابل انتخاب شده و وزن‌دهی می‌شود. این رویکرد سطح توکن، به الگوریتم اجازه می‌دهد تا شباهت‌های معنایی را بدون از دست دادن دقت موضعی و واژگانی شناسایی کند. خروجی این محاسبات به شکل سه شاخص دقت، بازخوانی و امتیاز \lr{F1} ارائه می‌شود. در سناریوی جعبه‌سیاه ما، مقدار متغیر هدف مستقیماً از طریق امتیاز \lr{F1} این معیار فرمول‌بندی می‌گردد تا تعادل پایدار میان پوشش محتوا و دقت واژگانی برقرار شود:
\begin{equation}\label{eq:app_lime_blackbox_target}
    r_k = \text{\lr{BERTScore-F1}}\left( \tilde{\mathbf{y}}^{(k)}, y_t \right)
\end{equation}
این تطبیق ساختاری، امکان اجرای اصولی خط‌پایه‌ی \lr{LIME} را در مسدودترین محیط‌های جعبه‌سیاه، بدون نیاز به دستکاری توابع ارزیاب نهایی پژوهش، با موفقیت فراهم آورده است.
